% !PRINTED 542
\weberpage{542}
%<<CONT
Modul $\mu$, wenn die ganzen Zahlen $\xi, \xi'$ congruent sind. Hieraus
ergiebt sich, dass, wenn $\xi$ ein volles Restsystem nach dem Modul
$\mu$ durchläuft, dasselbe auch von dem Producte $\alpha\xi$ gilt,
wodurch der Satz bewiesen ist:

1. \gesperrt{Ist $\alpha$ eine ganze Zahl in $\Omega$, relativ prim zu
dem Modul $\mu$, ferner $\gamma$ eine beliebige ganze Zahl in $\Omega$,
so ist die Congruenz:}
\weq{1}{\alpha\xi \equiv \gamma \ (\mathrm{mod}\ \mu)}
\gesperrt{immer durch eine ganze Zahl $\xi$ lösbar, und auch nur durch
eine, wenn für $\xi$ ein volles Restsystem nach dem Modul $\mu$
vorgeschrieben ist.}

Wenn hierin $\mu$ selbst eine Zahl ist, die wir mit $\beta$ bezeichnen,
so ist auch der Quotient
\[
\frac{\alpha\xi - \gamma}{\beta} = -\eta
\]
eine ganze Zahl, und dann nimmt der vorstehende Satz die Form an:

2. \gesperrt{Sind $\alpha, \beta, \gamma$ drei ganze Zahlen in $\Omega$
und $\alpha, \beta$ relativ prim, so kann man zwei andere ganze Zahlen
$\xi, \eta$ in $\Omega$ so bestimmen, dass}
\weq{2}{\alpha\xi + \beta\eta = \gamma}
\gesperrt{wird. Insbesondere kann man also auch für zwei beliebige
relative Primzahlen $\alpha, \beta$ die Gleichung}
\weq{3}{\alpha\xi + \beta\eta = 1}
\gesperrt{durch ganze Zahlen $\xi, \eta$ befriedigen.}

Dieser Satz lässt sich auf ein System von mehreren Zahlen übertragen.

Wenn die Zahlen $\alpha, \beta, \gamma, \ldots$ in $\mathfrak{o}$ keinen
gemeinsamen Theiler haben, so können wir zunächst Zahlen $\eta_1, \zeta_1,
\ldots$ in $\mathfrak{o}$ so bestimmen, dass
\[
\beta_1 = \beta\eta_1 + \gamma\zeta_1 + \cdots
\]
relativ prim zu $\alpha$ wird. Wir haben nämlich, wenn $\pi_1, \pi_2,
\ldots$ die verschiedenen Primfactoren von $\alpha$ sind, deren keiner
in allen $\beta, \gamma, \ldots$ aufgehen kann, und wenn etwa $\beta$
durch $\pi_1$ nicht theilbar ist, $\eta_1$ durch $\pi_1$ untheilbar, die
übrigen Zahlen $\zeta_1, \ldots$ durch $\pi_1$ theilbar u.\,s.\,f.
anzunehmen, und erhalten für jede der Zahlen $\eta_1, \zeta_1, \ldots$ die
Bedingung, dass sie durch einige der Primfactoren
%CONT>>
